\documentclass[10pt,a4paper]{amsart}
\usepackage[margin=19mm]{geometry}
\usepackage{amsmath,amssymb,mathtools}
\usepackage[scr=rsfs,cal=euler]{mathalfa}
\usepackage{xfrac}
\usepackage[unicode,hidelinks,bookmarks=false]{hyperref}
\newcommand{\ie}{i.e.\ }
\newcommand{\cf}{cf.\ }
\newcommand{\resp}{respectively\ }
\newcommand{\Subsection}{\subsection}
\providecommand{\nobreakdash}{\nobreak\hbox{-}}
\setlength{\emergencystretch}{2em}
\multlinegap0pt
\usepackage{amssymb}
\usepackage[shortlabels]{enumitem}
\usepackage{xcolor}
\usepackage{tikz}
\newtheorem{proposition}{Proposition}
\usepackage{cleveref}
\crefname{proposition}{Proposition}{Propositions}
\newcommand{\NN}{\mathbb N}
\newcommand{\cQ}{\mathcal Q}
\title{Regularity of Symbolic Powers of Co-Chordal Edge Ideals}
\author{Chwas Ahmed, Mohammed Rafiq Namiq}
\date{Source: arXiv:2608.20542v1 (CC BY 4.0)}
\begin{document}
\maketitle

\noindent\emph{Source and licence.} Regularity of Symbolic Powers of Co-Chordal Edge Ideals. Version arXiv:2608.20542v1 (CC BY 4.0), distributed by arXiv under the Creative Commons Attribution 4.0 International licence, http://creativecommons.org/licenses/by/4.0/, as recorded in the arXiv licence metadata for this exact version and shown on the abstract page https://arxiv.org/abs/2608.20542v1. Full licence notice: \texttt{licenses/arxiv-2608.20542-CC-BY-4.0.txt}.\par\smallskip

\noindent\textit{Editorial scope.} Counted unit: the article's proposition proposition (source0.tex lines 519--521). The article's complete original proof of it is delivered with it (source0.tex lines 523--532, one \texttt{proof} environment of 10 lines). No mathematics is written, repaired or re-derived: the statement and the proof are the article's own lines, in the article's own notation, and no retained line is substituted or re-flowed.  Retained uncounted as the source-local context the proof needs: lines 120--122.  Nothing was omitted from any retained span, so the package carries no omission record.  No retained line contains a citation, so the wrapper carries no bibliography and no bibliography entry.  No retained line contains a figure environment, an image-inclusion command or drawing code, so no image is delivered and the package declares no asset dependency.  Editorial formatting only, in addition: an AMS \texttt{amsart} preamble that restates the article's own declarations and macros used by the retained lines, namely the environments \texttt{proposition} on separate counters, together with the standard packages those lines need.  The retained environments print in this document as Proposition 1 in order of appearance, whereas the article prints the same items with the numbering of its own full document.  The complete native source of the article as distributed in the arXiv e-print for this exact version is bundled unchanged as \texttt{sources/208183/source0.tex}.
\begin{equation}\label{eq:1}
I(G)^{(s)}=\bigcap_{Q\in\cQ(H)}P_{C_Q}^s.
\end{equation}

\begin{proposition}\label[proposition]{prop:3}
Let $G$ be co-chordal, let $s\ge1$, and put $J=I(G)^{(s)}$. The ideal $J_{\langle s+1\rangle}$ is either zero or has linear quotients. In particular, whenever it is nonzero, it has an $(s+1)$-linear resolution.
\end{proposition}

\begin{proof}
Let $H=G^c$ and let $\mathbf a\in\NN^n$ satisfy $|\mathbf a|=s+1$. By \cref{eq:1}, $x^{\mathbf a}\in J$ if and only if $\sum_{i\notin Q}a_i\ge s$ for every $Q\in\cQ(H)$. Since $|\mathbf a|=s+1$, this is equivalent to $\mathbf a(Q)\le1$ for every maximal clique $Q$ of $H$. Every vertex belongs to a maximal clique, so $a_i\le1$ for every $i$. Thus every monomial of $J_{\langle s+1\rangle}$ is squarefree, and its support meets each clique of $H$ in at most one vertex. Equivalently,
\begin{equation}\label{eq:9}
J_{\langle s+1\rangle}
=
\bigl(x_A:A\subseteq[n],\ |A|=s+1,\ A\text{ is independent in }H\bigr).
\end{equation}

Choose a perfect elimination ordering $1,\ldots,n$ of $H$ and order the generators in \cref{eq:9} decreasingly in the lexicographic order induced by $x_1>\cdots>x_n$. Let $x_A>x_B$ be two generators, and let $i$ be the first index at which the characteristic vectors of $A$ and $B$ differ. Then $i\in A\setminus B$. If $k<i$ and $k\in B$, then $k\in A$, so $k$ is not adjacent to $i$ in $H$. The later neighbors of $i$ form a clique, while $B$ is independent, so $B$ contains at most one later neighbor of $i$. If such a neighbor exists, call it $j$. It necessarily belongs to $B\setminus A$. If it does not exist, choose any $j\in B\setminus A$. In either case $C=(B\setminus\{j\})\cup\{i\}$ is independent in $H$. Also $j>i$, so $x_C>x_B$, $(x_C:x_B)=x_i$, and $x_i\mid(x_A:x_B)$. This is the linear quotient criterion for the above lexicographic ordering. Since all generators have degree $s+1$, the resolution is $(s+1)$-linear.
\end{proof}

\end{document}
